\section{Limitations}
\label{sec:limitations}

Our study is bounded in four ways. \textbf{(1) Stack maturity.} The AMD results use
a new accelerator (CDNA4 / gfx950) on a young ROCm software stack; several
configurations did not run and are reported as explicit coverage gaps rather than
omitted---TGI does not serve gfx950, the vLLM fused-MoE MXFP4 loader cannot read the
35B checkpoint, the 235B's dual-chunk 1M context does not load, and HIP graphs
regress throughput---so the AMD numbers reflect today's stack and should improve as
it matures. \textbf{(2) Cross-vendor scope.} The NVIDIA B200 comparison is
single-node (one node available) and uses the 235B only; its power telemetry was
not collected, so B200 energy and tokens/joule are unavailable, and the comparison
assumes each platform at its best-performing configuration (eager on AMD).
\textbf{(3) Models and modality.} We evaluate two checkpoints; the medium tier is a
multimodal MoE served text-only, so its vision encoder is loaded but unexercised
(and is the cause of the FP8/TP$=$1 constraint). \textbf{(4) Quality guardrail.} The
guardrail is a lightweight exact-match needle-recall check designed to catch
speed-for-quality regressions, not a full task-quality evaluation; absolute
correctness on real RAG tasks is out of scope. Finally, several capacity points sit
at the first ladder rung, where goodput is small and prefill-dominated; we report
run-to-run variance for these boundary points.
